\chapter{Design}

\section{A primary I/O class for each source/sink type}
The library divides I/O functionality into a primary I/O class
for a given type of source and sink.  For example,
\hyperref[section:DwmStreamIOClass]{\texttt{Dwm::StreamIO}} provides
functions for \texttt{std::istream} sources and \texttt{std::ostream}
sinks.  Table~\ref{table:PrimaryClassSourceSinkType} lists the primary
I/O classes and their associated source and sink type.

\begin{table}[H]
\begin{tabular}{llll}
  \textbf{Primary I/O class} & \textbf{Source type} & \textbf{Sink type} & \textbf{Description} \\ \hline
  \hyperref[section:DwmStreamIOClass]{\texttt{Dwm::StreamIO}} & \texttt{std::istream \&} & \texttt{std::ostream \&} & references to C++ iostreams \\
  \hyperref[section:DwmFileIOClass]{\texttt{Dwm::FileIO}} & \texttt{FILE *} & \texttt{FILE *} & \texttt{FILE} created with \texttt{fopen()} and friends \\
  \hyperref[section:DwmDescriptorIOClass]{\texttt{Dwm::DescriptorIO}} & \texttt{int} & \texttt{int} & UNIX descriptors \\
  \hyperref[section:DwmBZ2IOClass]{\texttt{Dwm::BZ2IO}} & \texttt{BZFILE *} & \texttt{BZFILE *} & \texttt{BZFILE} created with \texttt{BZ2\_bzopen()}\\
  \hyperref[section:DwmGZIOClass]{\texttt{Dwm::GZIO}} & \texttt{gzFile} & \texttt{gzFile} & \texttt{gzFile} created with \texttt{gzopen()} \\
\end{tabular}
\caption{Primary I/O class source and sink types}
\label{table:PrimaryClassSourceSinkType}
\end{table}

Each of the primary I/O classes provides a set of static functions and
function templates to read and write data.

\section{Conceptual interfaces}
Conceptually, the I/O classes provide four primary interfaces to read
and write from/to their associated source/sink type, as depicted in
table~\ref{table:ConceptualInterfaceGeneric}.  There are
single-object interfaces (\texttt{ReadFunc} and \texttt{WriteFunc})
and multi-object variadic interfaces (\texttt{VariadicReadFunc} and
\texttt{VariadicWriteFunc}).  The multi-object variadic interfaces
just call to the single-object interfaces for each object.

\begin{table}[H]
\begin{tabular}{l}
  \textbf{Read/write functions (conceptual)} \\ \hline
  \makecell[l]{\texttt{template <typename T>}\\ \texttt{static \textit{\color{red1}ReadReturnType} \& \color{darkblue}{\textit{ReadFunc}}(\textit{\color{green1}SourceType} {\color{green1}source}, T \& t);}} \\
  \makecell[l]{\texttt{template <typename T>}\\ \texttt{static \textit{\color{red1}WriteReturnType} \textit{\color{darkblue}WriteFunc}(\textit{\color{green1}SinkType} {\color{green1}sink}, const T \& t);}} \\
  \makecell[l]{\texttt{template <typename... Args>}\\ \texttt{static \textit{\color{red1}ReadReturnType} \textit{\color{darkblue}VariadicReadFunc}(\textit{\color{green1}SourceType} {\color{green1}source}, Args \& ...args);}} \\
  \makecell[l]{\texttt{template <typename... Args>}\\ \texttt{static \textit{\color{red1}WriteReturnType} \textit{\color{darkblue}VariadicWriteFunc}(\textit{\color{green1}SinkType} {\color{green1}sink}, const Args \& ...args);}} \\
\end{tabular}
\caption{Conceptual read/write functions in all primary I/O classes}
\label{table:ConceptualInterfaceGeneric}
\end{table}

Table~\ref{table:ConceptualInterfaceMapping} maps the generic conceptual
interface to each of the primary I/O classes.

\begin{table}[H]
\begin{tabular}{l|l|l|l|l|l}
  & \multicolumn{5}{c}{\texttt{Dwm} I/O class} \\
  & \texttt{StreamIO} & \texttt{FileIO} & \texttt{DescriptorIO} & \texttt{BZ2IO} & \texttt{GZIO} \\ \hline
  \texttt{\textit{\color{darkblue}ReadFunc}} & \texttt{\color{darkblue}Read} & \texttt{\color{darkblue}Read} & \texttt{\color{darkblue}Read} & \texttt{\color{darkblue}BZRead} & \texttt{\color{darkblue}Read} \\
  \texttt{\textit{\color{darkblue}VariadicReadFunc}} & \texttt{\color{darkblue}ReadV} & \texttt{\color{darkblue}ReadV} & \texttt{\color{darkblue}ReadV} & \texttt{\color{darkblue}BZReadV} & \texttt{\color{darkblue}ReadV} \\
  \texttt{\textit{\color{red1}ReadReturnType}} & \texttt{\color{red1}std::istream \&} & \texttt{\color{red1}size\_t} & \texttt{\color{red1}ssize\_t} & \texttt{\color{red1}int} & \texttt{\color{red1}int} \\
  \color{green1}\texttt{\textit{SourceType}} & \color{green1}\texttt{std::istream \&} & \color{green1}\texttt{FILE *} & \color{green1}\texttt{int} & \color{green1}\texttt{BZFILE *} & \color{green1}\texttt{gzFile} \\
  \texttt{\textit{\color{darkblue}WriteFunc}} & \texttt{\color{darkblue}Write} & \texttt{\color{darkblue}Write} & \texttt{\color{darkblue}Write} & \texttt{\color{darkblue}BZWrite} & \texttt{\color{darkblue}Write} \\
  \texttt{\textit{\color{darkblue}VariadicWriteFunc}} & \texttt{\color{darkblue}WriteV} & \texttt{\color{darkblue}WriteV} & \texttt{\color{darkblue}WriteV} & \texttt{\color{darkblue}BZWriteV} & \texttt{\color{darkblue}WriteV} \\
  \texttt{\textit{\color{red1}WriteReturnType}} & \texttt{\color{red1}std::ostream \&} & \texttt{\color{red1}size\_t} & \texttt{\color{red1}ssize\_t} & \texttt{\color{red1}int} & \texttt{\color{red1}int} \\
  \texttt{\textit{\color{green1}SinkType}} & \texttt{\color{green1}std::ostream \&} & \texttt{\color{green1}FILE *} & \texttt{\color{green1}int} & \texttt{\color{green1}BZFILE *} & \texttt{\color{green1}gzFile} \\
\end{tabular}
\caption{Conceptual interface mapping to the primary I/O classes}
\label{table:ConceptualInterfaceMapping}
\end{table}

\section{Categories of interface overloads in actual implementation}
There are three categories of overloads in the implementation:
\begin{itemize}[noitemsep,nolistsep]
\item Regular functions for directly supported types
\item Function templates for STL containers and other standard class
  templates
\item Function templates for types that meet the conceptual requirements
  of a given \hyperref[chapter:primaryIOClasses]{primary I/0 class}
\end{itemize}

\subsection{Regular functions for directly supported types}
Reading and writing of directly supported types is done with regular
function overloads.  For example, this function from
\hyperref[section:DwmStreamIOClass]{\texttt{Dwm::StreamIO}}
to write a \texttt{uint32\_t} value:
\begin{verbatim}
  static std::ostream & Write(std::ostream & os, uint32_t value);
\end{verbatim}

\subsection{Function templates for STL containers and other standard class templates}
Overloaded function templates are used to read and write instantiations of
standard class templates, where the parameter types of the standard
class templates are either directly supported types, types that meet
the conceptual requirements of a readable and writable class, or instances
of supported standard templates using any of the aforementioned types in
type parameters.  In other words, supported standard containers of any depth
can be read and written, as long as the leaf types are directly
supported types or types that meet the conceptual requirements of a
readable and writable class for the I/O class of interest.

This is probably best explained via some example readable and writable
types:
\begin{verbatim}
  std::vector<std::string>
  std::map<uint32_t,std::vector<std::string>>
  std::map<uint32_t,std::vector<std::map<uint8_t,std::string>>>
\end{verbatim}

An example of one such function template from
\hyperref[section:DwmStreamIOClass]{\texttt{Dwm::StreamIO}}:

\begin{verbatim}
    template <typename valueT, size_t N>
    static std::istream & Read(std::istream & is, std::array<valueT, N> & a)
\end{verbatim}

\subsection{Function templates for types that meet the conceptual requirements of a given primary I/O class}
Overloaded function templates with concept requirements are used to read
and write types that meet the concept requirements.  The concepts required
for each primary I/O class are listed in
table~\ref{table:PrimaryIOClassConcepts}.

An example of one such function template from
\hyperref[section:DwmStreamIOClass]{\texttt{Dwm::StreamIO}}:

\begin{ttfamily}
  static std::ostream \& Write(std::ostream \& os, const \hyperlink{concept:HasStreamWrite}{HasStreamWrite} auto \& t)
\end{ttfamily}

\begin{table}[H]
\begin{tabular}{lll}
  \textbf{\texttt{Dwm} I/O class} & \textbf{Read concept} & \textbf{Write concept} \\ \hline
  \hyperref[section:DwmStreamIOClass]{\texttt{StreamIO}} & \hyperlink{concept:HasStreamRead}{\texttt{Dwm::HasStreamRead<T>}} & \hyperlink{concept:HasStreamWrite}{\texttt{Dwm::HasStreamWrite<T>}} \\
  \hyperref[section:DwmFileIOClass]{\texttt{FileIO}} & \hyperlink{concept:HasFileRead}{\texttt{Dwm::HasFileRead<T>}} & \hyperlink{concept:HasFileWrite}{\texttt{Dwm::HasFileWrite</t>}} \\
  \hyperref[section:DwmDescriptorIOClass]{\texttt{DescriptorIO}} & \hyperlink{concept:HasDescriptorRead}{\texttt{Dwm::HasDescriptorRead<T>}} & \hyperlink{concept:HasDescriptorWrite}{\texttt{Dwm::HasDescriptorWrite<T>}} \\
  \hyperref[section:DwmBZ2IOClass]{\texttt{BZ2IO}} & \hyperlink{concept:HasBZRead}{\texttt{Dwm::HasBZRead<T>}} & \hyperlink{concept:HasBZWrite}{\texttt{Dwm::HasBZWrite<T>}} \\
  \hyperref[section:DwmGZIOClass]{\texttt{GZIO}} & \hyperlink{concept:DwmHasGZRead}{\texttt{Dwm::HasGZRead<T>}} & \hyperlink{concept:DwmHasGZWrite}{\texttt{Dwm::HasGZWrite<T>}} \\
\end{tabular}
\caption{Primary I/O class concepts for user-defined types}
\label{table:PrimaryIOClassConcepts}
\end{table}

\section{\texttt{Dwm::StreamIO} read/write functions}
\subsection{\texttt{Dwm::StreamIO} conceptual read/write functions}
Table ~\ref{table:StreamIOConceptualReadWriteFuncs} lists the conceptual
read and write member functions in \texttt{Dwm::StreamIO}.
\begin{table}[H]
\begin{tabular}{ll}
  \textbf{Primary I/O class} & \textbf{Read/write functions (conceptual)} \\ \hline
  \multirow{4}{*}{\hyperref[section:DwmStreamIOClass]{\texttt{Dwm::StreamIO}}} & \makecell[l]{\texttt{template <typename T>}\\ \texttt{static std::istream \& Read(std::istream \& is, T \& t);}} \\ \cline{2-2}
  & \makecell[l]{\texttt{template <typename T>}\\ \texttt{static std::ostream \& Write(std::ostream \& os, const T \& t);}} \\ \cline{2-2}
  & \makecell[l]{\texttt{template <typename... Args>}\\ \texttt{static std::istream \& ReadV(std::istream \& is, Args \& ...args);}} \\ \cline{2-2}
  & \makecell[l]{\texttt{template <typename... Args>}\\ \texttt{static std::ostream \& WriteV(std::ostream \& os, const Args \& ...args);}} \\ \hline
\end{tabular}
\caption{Conceptual read and write member functions in \texttt{Dwm::StreamIO}}
\label{table:StreamIOConceptualReadWriteFuncs}
\end{table}

\begin{minipage}{\textwidth}
Conceptually...

\subsubsection{\texttt{template <typename T>}\\ \texttt{static std::istream \& Read(std::istream \& is, T \& t)}}
\begin{itemize}[noitemsep,nolistsep]
\item Reads the contents of \texttt{t} from \texttt{is} and returns
  \texttt{is}.
\end{itemize}

\subsubsection{\texttt{template <typename T>}\\ \texttt{static std::Write(std::ostream \& os, const T \& t)}}
\begin{itemize}[noitemsep,nolistsep]
\item Writes the contents of \texttt{t} to \texttt{os} and returns
  \texttt{os}.
\end{itemize}

\subsubsection{\texttt{template <typename... Args>}\\ \texttt{static std::istream \& ReadV(std::istream \& is, Args \& ...args)}}
\begin{itemize}[noitemsep,nolistsep]
\item Reads the contents of all \texttt{args} (in order) from \texttt{is}
  and returns \texttt{is}.
\end{itemize}

\subsubsection{\texttt{template <typename... Args>}\\ \texttt{static std::ostream \& WriteV(std::ostream \& os, const Args \& ...args)}}
\begin{itemize}[noitemsep,nolistsep]
\item Writes the contents of all \texttt{args} (in order) to \texttt{os}
  and returns \texttt{os}.
\end{itemize}

\end{minipage}

\section{\texttt{Dwm::FileIO} read/write functions}
\subsection{\texttt{Dwm::FileIO} conceptual read/write functions}
\begin{tabular}{ll}
  \textbf{Primary I/O class} & \textbf{Read/write functions (conceptual)} \\ \hline  \multirow{4}{*}{\hyperref[section:DwmFileIOClass]{\texttt{Dwm::FileIO}}} & \makecell[l]{\texttt{template <typename T>}\\ \texttt{static size\_t Read(FILE *f, T \& t);}} \\ \cline{2-2}
  & \makecell[l]{\texttt{template <typename T>}\\ \texttt{static size\_t Write(FILE *f, const T \& t);}} \\ \cline{2-2}
  & \makecell[l]{\texttt{template <typename... Args>}\\ \texttt{static size\_t ReadV(FILE *f, Args \& ...args);}} \\ \cline{2-2}
  & \makecell[l]{\texttt{template <typename... Args>}\\ \texttt{static size\_t WriteV(FILE *f, const Args \& ...args);}} \\ \hline
\end{tabular}

\section{\texttt{Dwm::DescriptorIO} read/write functions}
\subsection{\texttt{Dwm::DescriptorIO} conceptual read/write functions}
\begin{tabular}{ll}
  \textbf{Primary I/O class} & \textbf{Read/write functions (conceptual)} \\ \hline
  \multirow{2}{*}{\hyperref[section:DwmDescriptorIOClass]{\texttt{Dwm::DescriptorIO}}} & \makecell[l]{\texttt{template <typename T>}\\ \texttt{static ssize\_t Read(int fd, T \& t);}} \\ \cline{2-2}
  & \makecell[l]{\texttt{template <typename T>}\\\texttt{static ssize\_t Write(int fd, const T \& t);}} \\ \cline{2-2}
  & \makecell[l]{\texttt{template <typename... Args>}\\ \texttt{static ssize\_t ReadV(int fd, Args \& ...args);}} \\ \cline{2-2}
  & \makecell[l]{\texttt{template <typename... Args>}\\ \texttt{static ssize\_t WriteV(int fd, const Args \& ...args);}} \\ \hline
\end{tabular}

\section{\texttt{Dwm::BZ2IO} read/write functions}
\subsection{\texttt{Dwm::BZ2IO} conceptual read/write functions}
\begin{tabular}{ll}
  \textbf{Primary I/O class} & \textbf{Read/write functions (conceptual)} \\ \hline
  \multirow{2}{*}{\hyperref[section:DwmBZ2IOClass]{\texttt{Dwm::BZ2IO}}} & \makecell[l]{\texttt{template <typename T>}\\ \texttt{static int BZRead(BZFILE *bzf, T \& t);}} \\ \cline{2-2}
  & \makecell[l]{\texttt{template <typename T>}\\ \texttt{static int BZWrite(BZFILE *bzf, const T \& t)}} \\ \cline{2-2}
 & \makecell[l]{\texttt{template <typename... Args>}\\ \texttt{static int BZReadV(BZFILE *bzf, Args \& ...args);}} \\ \cline{2-2}
  & \makecell[l]{\texttt{template <typename T>}\\ \texttt{static int BZWriteV(BZFILE *bzf, const Args \& ...args)}} \\ \hline
\end{tabular}

\section{\texttt{Dwm::GZIO} read/write functions}
\subsection{\texttt{Dwm::GZIO} conceptual read/write functions}
\begin{tabular}{ll}
  \textbf{Primary I/O class} & \textbf{Read/write functions (conceptual)} \\ \hline
  \multirow{2}{*}{\hyperref[section:DwmGZIOClass]{\texttt{Dwm::GZIO}}} & \makecell[l]{\texttt{template <typename T>}\\ \texttt{static int Read(gzFile gzf, T \& t);}} \\ \cline{2-2}
  & \makecell[l]{\texttt{template <typename T>}\\ \texttt{static int Write(gzFile gzf, const T \& t);}} \\ \cline{2-2}
 & \makecell[l]{\texttt{template <typename... Args>}\\ \texttt{static int ReadV(gzFile gzf, Args \& ...args);}} \\ \cline{2-2}
  & \makecell[l]{\texttt{template <typename... Args>}\\ \texttt{static int Write(gzFile gzf, const Args \& ...args);}} \\ \hline
\end{tabular}

\section{Member function overloads}
libDwm uses member function overloading to allow the user to use the
same function names to read or write different types from/to a
source/sink.  This is also critical to supporting instances of class
templates.

\section{Member function templates}
libDwm uses function templates and function template overloads to
provide the same interfaces (the same member function names) to
read or write instances of standard class templates from/to a
source/sink (\texttt{std::vector}, \texttt{std::pair},
\texttt{std::tuple}, \texttt{std::map}, et. al.).

\section{Concepts}
libDwm uses concepts in many places to aid in overload resolution
and to document the requirements of user-defined types.

\section{Support for user-designed types}
\subsection{Before C++26}
In general, code utilized with pre-C++26 compilers and standard libraries
requires a pair of member functions to be added to any user-defined type
for a given source/sink pair.  Much of the time, these functions are
fairly simple to implement, particuarly for C++ iostreams
(\texttt{std::istream} as a source, \texttt{std::ostream} as a sink).
This is especially true for user-defined types that only contain types
already supported by the corresponding primry I/O class (for C++ iostreams,
that would be \texttt{Dwm::StreamIO}).

\subsection{C++26 and beyond}
With C++26 features (particularly P2996 and P3394), it becomes possible
to elegantly handle serialization and deserialization of many user
defined types with no changes to user code.  It also allows annotations
of types and variables that can control the serialization and
deserialization of a user defined type (for example, skipping a data
member of a class during serialization and deserialization).  Work
here has already begun, using a slightly modified fork of the bloomberg
clang implementation.
